% Lettura: "Funziona a meraviglia", Wenke Husmann (Internazionale, 22 gennaio 2026)
\documentclass[11pt]{report}
% ============================================================
%  ENCODING E FONT
% ============================================================
\usepackage{fontspec}
\setmainfont{TeX Gyre Pagella}
% ============================================================
%  LINGUE
% ============================================================
\usepackage{polyglossia}
\setmainlanguage{italian}
\setotherlanguage[variant=ancient]{greek}

% ============================================================
%  MATEMATICA
% ============================================================
\usepackage{amsmath}
\usepackage{amssymb}
\usepackage{amsfonts}
\usepackage{amsthm}
\usepackage{mathtools}
\usepackage{physics}
\usepackage{yhmath}
\usepackage{cancel}
\usepackage{bbm}

% unicode-math va caricato dopo amsmath/amssymb/amsfonts (altrimenti conflitti
% tipo "\eth already defined"); fornisce il font matematico per XeLaTeX.
\usepackage{unicode-math}
\setmathfont{TeX Gyre Pagella Math}

%dimensionimatematica
\makeatletter
\DeclareMathSizes{\@xpt}{10}{7.5}{6}    % testo 10pt
\DeclareMathSizes{\@xipt}{10}{7.5}{6}   % testo 11pt
\DeclareMathSizes{\@xiipt}{10}{7.5}{6}  % testo 12pt
\makeatother

% Spaziatura display math
\setlength{\abovedisplayskip}{4pt}
\setlength{\belowdisplayskip}{4pt}
\setlength{\abovedisplayshortskip}{2pt}
\setlength{\belowdisplayshortskip}{2pt}

% Comandi matematici
\newcommand{\R}{\mathbb{R}}
\newcommand{\sse}{ \ \textup{sse} \ }
\newcommand{\tc}{ \textup{ t.c. }}
\newcommand{\under}[1]{_{\textup{#1}}}
\newcommand{\sopra}[1]{^{\textup{#1}}}
\newcommand{\tonde}[1]{\left(#1\right)}
\newcommand{\quadre}[1]{\left[#1\right]}
\newcommand{\graffe}[1]{\left\lbrace #1 \right\rbrace}
\newcommand{\et}{\wedge}
\newcommand{\vel}{\vee}
\newcommand{\aut}{\dot{\vee}}
\newcommand{\quindi}{\quad \Rightarrow \quad}
\newcommand{\ubar}[1]{\underline{\mathord{#1}}}

% ============================================================
%  UNITÀ DI MISURA
% ============================================================
\usepackage{siunitx}
\sisetup{
  detect-all,
  output-decimal-marker = {,},
  per-mode = symbol
}
\DeclareSIPrefix\micro{μ}{-6}
% ============================================================
%  LAYOUT E GEOMETRIA
% ============================================================
\usepackage[centering, includeheadfoot]{geometry}

\newlength{\margineEsterno}
\newlength{\margineInterno}
\newlength{\margineAlto}
\newlength{\margineBasso}

\setlength{\margineEsterno}{3.5cm}
\setlength{\margineInterno}{1cm}
\setlength{\margineAlto}{1cm}
\setlength{\margineBasso}{1cm}

\geometry{
  paperwidth  = 170mm,
  paperheight = 240mm,
  top         = \margineAlto,
  bottom      = \margineBasso,
  inner       = \margineInterno,
  outer       = \margineEsterno,
  twoside,
}

% Dimensioni margine per QR e note
\newlength{\maxQR}
\newlength{\sizeQR}
\newlength{\largQR}

\setlength{\maxQR}{2.5cm}
\setlength{\marginparsep}{3mm}
\setlength{\marginparwidth}{\dimexpr\margineEsterno - \marginparsep\relax}
\setlength{\largQR}{\dimexpr\marginparwidth - 6mm\relax}

\newcommand{\calcsizeQR}{%
  \setlength{\sizeQR}{\dimexpr\marginparwidth - 10pt\relax}%
  \ifdim\sizeQR>\maxQR \setlength{\sizeQR}{\maxQR}\fi
}

% Spaziatura testo
\usepackage{setspace}
\setstretch{1.18}
\usepackage{parskip}
\setlength{\parindent}{0pt}
\setlength{\parskip}{0.25\baselineskip plus 1pt minus 1pt}

%Sforare le immagini nel margine esterno
\newcommand{\sforaEsterno}[1]{%
  \checkoddpage
  \ifoddpage
    \begin{adjustwidth}{0pt}{-\margineEsterno}#1\end{adjustwidth}%
  \else
    \begin{adjustwidth}{-\margineEsterno}{0pt}#1\end{adjustwidth}%
  \fi
}

\raggedbottom

% ============================================================
%  INTESTAZIONI E PIÈ DI PAGINA
% ============================================================
\usepackage{fancyhdr}
\pagestyle{fancy}
\fancyhf{}
\fancyfoot[C]{\thepage}
\renewcommand{\headrulewidth}{0pt}
\renewcommand{\footrulewidth}{0pt}
\let\cleardoublepage\clearpage
\renewcommand{\cleardoublepage}{\clearpage}

% ============================================================
%  TITOLI E SEZIONI
% ============================================================
\usepackage{titlesec}
\usepackage{tocloft}

% Part
\renewcommand{\part}[1]{%
  \refstepcounter{part}
  \clearpage
  \thispagestyle{empty}
  \begin{center}
    \Huge\bfseries Part \thepart \\[1em] #1
  \end{center}
  \vspace{2em}
  \addcontentsline{toc}{part}{Part \thepart: #1}
}

% Helper: salva il nome del capitolo E lo stampa
\newcommand{\savechaptertitle}[1]{\gdef\currentchaptername{#1}#1}

\titleformat{\chapter}[block]
  {\normalfont\LARGE\bfseries\color{cyan}}
  {\rule{\textwidth}{.5pt}\\[0.1em]\centering\Roman{chapter}.\,}
  {1em}
  {\huge\scshape\centering\savechaptertitle}
  [\vspace{-0.2em}\rule{\textwidth}{0.5pt}]

% ============================================================
%  CAPITOLO CON IMMAGINE OPZIONALE
% ============================================================
\newcommand{\chapterimage}[3]{%
  \clearpage
  \thispagestyle{fancy}%
  \begin{adjustwidth}{-\margineEsterno}{0pt}%
    \includegraphics[width=\dimexpr\textwidth+\margineEsterno\relax]{#1}%
    \par\vspace{2pt}%
    {\small\itshape\color{gray} #2}%
  \end{adjustwidth}%
  \vspace{0.5em}%
  \refstepcounter{chapter}%
  \gdef\currentchaptername{#3}%
  \addcontentsline{toc}{chapter}{\protect\numberline{\thechapter}#3}%
  \begin{center}
    {\color{cyan}\rule{\textwidth}{.5pt}}\\[0.2em]
    {\color{cyan}\LARGE\bfseries\Roman{chapter}.\,}\\[0.1em]
    {\color{cyan}\huge\scshape #3}\\[-0.3em]
    {\color{cyan}\rule{\textwidth}{.5pt}}
  \end{center}
  \vspace{0.5em}%
}

\titleformat{\section}[block]
  {\normalfont\Large\bfseries\color{cyan}}{\thesection}{1em}{}

\titleformat{\subsection}[block]
  {\normalfont\large\bfseries\color{cyan}}{\thesubsection}{1em}{}

\titleformat{\subsubsection}[block]
  {\normalfont\bfseries\color{cyan}}{\thesubsubsection}{1em}{}

\titleformat{\paragraph}[block]
  {\normalfont\bfseries\color{cyan}}{}{0em}{}

\titlespacing*{\chapter}   {0pt}{-30pt}{8pt}
\titlespacing*{\section}   {0pt}{0.8\baselineskip}{0.3\baselineskip}
\titlespacing*{\subsection}{0pt}{0.6\baselineskip}{0.2\baselineskip}
\setlength{\cftbeforechapskip}{4pt}

% ============================================================
%  LISTE
% ============================================================
\usepackage[inline]{enumitem}
\setlist[itemize,1]{label=\textcolor{gray!70}{\rule{0.8ex}{0.8ex}}}
\setlist{itemsep=0pt, parsep=0pt, topsep=3pt, partopsep=0pt}
\setlength{\topsep}{2pt}
\newlist{myenum}{enumerate*}{1}
\setlist[myenum]{label=\arabic*, itemjoin={\hspace{1.5cm}}}
\newlist{sceltamultipla}{enumerate}{1}
\setlist[sceltamultipla,1]{%
  label=\textcolor{gray}{\fbox{\textbf{\footnotesize\makebox[1.2em]{\Alph*}}}},
  leftmargin=*,
  itemsep=1pt,
  topsep=2pt
}

% ============================================================
%  COLORI
% ============================================================
\usepackage{xcolor}

\definecolor{ciano}           {RGB}{0,   180, 200}
\definecolor{sfondoImportante}{RGB}{255, 252, 220}
\definecolor{coloreAnimazione}{RGB}{0,   150, 136}
\definecolor{sfondoAnimazione}{RGB}{225, 245, 243}
\definecolor{gold}            {RGB}{255, 215,   0}
\definecolor{copper}          {RGB}{104, 115,  51}

\definecolor{buildSub}  {HTML}{FFF0F5}  \definecolor{buildTitle} {HTML}{C71585}
\definecolor{readSub}   {HTML}{F2F2F2}  \definecolor{readTitle}  {HTML}{333333}
\definecolor{deriveSub} {HTML}{FFF9E6}  \definecolor{deriveTitle}{HTML}{B8860B}
\definecolor{exSub}     {HTML}{F0F7FF}  \definecolor{exTitle}    {HTML}{005C99}
\definecolor{impSub}    {HTML}{FFF0F0}  \definecolor{impTitle}   {HTML}{CC0000}
\definecolor{warnSub}   {HTML}{FFF0F0}  \definecolor{warnTitle}  {HTML}{B22222}

\definecolor{sfondoRiassunto}{RGB}{215, 242, 240}
\definecolor{intestRiassunto}{RGB}{  0, 145, 132}
\definecolor{rsezColor}      {RGB}{  0, 160, 146}

% ============================================================
%  TIKZ E GRAFICI
% ============================================================
\usepackage{tikz}
\usepackage{tikz-3dplot}
\usepackage{pgfplots}
\usepackage{circuitikz}
\usepackage{tikzsymbols, fontawesome5}

\usetikzlibrary{
  decorations.pathreplacing, decorations.markings, decorations.pathmorphing,
  angles, quotes, arrows.meta, calc, positioning,
  shapes, shapes.symbols, shadows, fit,
  backgrounds, matrix, patterns, tikzmark
}

\pgfplotsset{compat=1.17}
\usepgfplotslibrary{colormaps}
\tdplotsetmaincoords{70}{110}

% Simbolo caution
\newcommand{\caution}[1]{%
  \begingroup
  \pgfmathsetmacro{\s}{#1}%
  \pgfmathsetmacro{\Rr}{1.6*\s}%
  \pgfmathsetmacro{\Lwd}{2*\s}%
  \pgfmathsetmacro{\FSpt}{40*\s}%
  \pgfmathsetmacro{\BLSpt}{1.2*\FSpt}%
  \tikz[baseline=(bang.base)]{
    \draw[line width=\Lwd mm, red, fill=yellow!30]
      (90:\Rr) -- (210:\Rr) -- (330:\Rr) -- cycle;
    \node (bang) at (0,0)
      {\fontsize{\FSpt pt}{\BLSpt pt}\selectfont\bfseries !};
  }%
  \endgroup
}

% ============================================================
%  BOX: TCOLORBOX
% ============================================================
\usepackage[many]{tcolorbox}

\tcbset{
  mystyle/.style={
    enhanced, breakable,
    boxrule=0pt, frame hidden, arc=2mm,
    top=4mm, bottom=4mm, left=4mm, right=4mm,
    fonttitle=\bfseries\sffamily, coltitle=white,
    attach boxed title to top left={yshift=-2mm, xshift=0mm},
    boxed title style={
      empty,
      top=1.5mm, bottom=1.5mm, left=3mm, right=8mm,
      underlay={
        \fill[#1, rounded corners=2mm]
          (frame.south west) -- (frame.north west) --
          ([xshift=-5mm]frame.north east) -- (frame.south east) -- cycle;
      }
    }
  }
}

\newtcolorbox{mdwarning}[1][]{
  mystyle={warnTitle}, colback=warnSub,
  title={\ifstrempty{#1}{Attenzione}{#1}}
}
\newtcolorbox{mdbuild}[1][]{
  mystyle={buildTitle}, colback=buildSub,
  title={\ifstrempty{#1}{Costruire la formula}{#1}}
}
\newtcolorbox{mdread}[1][]{
  mystyle={readTitle}, colback=readSub,
  title={\ifstrempty{#1}{Leggere la formula}{#1}}
}
\newtcolorbox{mdderive}[1][]{
  mystyle={deriveTitle}, colback=deriveSub,
  title={\ifstrempty{#1}{Derivare la formula}{#1}}
}
\newtcolorbox{mdexample}[1][]{
  mystyle={exTitle}, colback=exSub,
  title={\ifstrempty{#1}{Esempio guidato}{#1}}
}

\newtcolorbox{mdimportant}[2][]{%
  enhanced, breakable,
  boxrule=0pt, frame hidden,
  colback=sfondoImportante, arc=1mm,
  top=4mm, bottom=4mm, left=6mm, right=4mm,
  overlay={
    \checkoddpage
    \ifoddpage
      \draw[ciano, line width=2pt] (frame.south east)--(frame.north east);
      \draw[ciano, line width=2pt] (frame.east)--([xshift=5mm]frame.east);
      \draw[ciano, fill=white, line width=2pt] ([xshift=5mm]frame.east) circle (4pt);
      \node[anchor=west, text=ciano, font=\bfseries\sffamily\small,
            align=left, text width=\dimexpr\largQR-2mm\relax]
        at ([xshift=7mm]frame.east) {#2};
    \else
      \draw[ciano, line width=2pt] (frame.south west)--(frame.north west);
      \draw[ciano, line width=2pt] ([xshift=-5mm]frame.west)--(frame.west);
      \draw[ciano, fill=white, line width=2pt] ([xshift=-5mm]frame.west) circle (4pt);
      \node[anchor=east, text=ciano, font=\bfseries\sffamily\small,
            align=right, text width=\dimexpr\largQR-2mm\relax]
        at ([xshift=-7mm]frame.west) {#2};
    \fi
  },
  #1
}

\newsavebox{\mdanimationQRbox}

\newtcolorbox{mdanimation}[2][Animazione]{%
  enhanced, breakable,
  boxrule=0pt, frame hidden,
  colback=sfondoAnimazione, arc=2mm,
  top=4mm, bottom=4mm, left=4mm, right=4mm,
  fonttitle=\bfseries\sffamily,
  coltitle=white,
  title={#1},
  attach boxed title to top left={yshift=-2mm, xshift=0mm},
  boxed title style={
    empty,
    top=1.5mm, bottom=1.5mm, left=3mm, right=8mm,
    underlay={
      \fill[coloreAnimazione, rounded corners=2mm]
        (frame.south west) -- (frame.north west) --
        ([xshift=-5mm]frame.north east) -- (frame.south east) -- cycle;
    }
  },
  overlay={
    \ifx\relax#2\relax\else
      \checkoddpage
      \calcsizeQR
      \ifoddpage
        \fill[sfondoAnimazione]
          ([xshift=-1mm, yshift=-2mm]frame.north east) --
          ([xshift=3mm, yshift=1mm]frame.north east) --
          ([xshift=3mm, yshift=-\the\sizeQR-0pt]frame.north east) --
          ([xshift=-1mm, yshift=-\the\sizeQR+3mm]frame.north east) -- cycle;
        \node[anchor=north west, draw=coloreAnimazione, line width=1.5pt,
              fill=sfondoAnimazione, inner sep=4pt, rounded corners=1mm]
          at ([xshift=3mm, yshift=2mm]frame.north east) {%
            \begin{minipage}{\the\largQR}\centering
              \qrcode[height=\the\sizeQR]{#2}
            \end{minipage}%
          };
      \else
        \fill[sfondoAnimazione]
          ([xshift=1mm, yshift=-2mm]frame.north west) --
          ([xshift=-3mm, yshift=1mm]frame.north west) --
          ([xshift=-3mm, yshift=-\the\sizeQR-0pt]frame.north west) --
          ([xshift=1mm, yshift=-\the\sizeQR+3mm]frame.north west) -- cycle;
        \node[anchor=north east, draw=coloreAnimazione, line width=1.5pt,
              fill=sfondoAnimazione, inner sep=4pt, rounded corners=1mm]
          at ([xshift=-3mm, yshift=2mm]frame.north west) {%
            \begin{minipage}{\the\largQR}\centering
              \qrcode[height=\the\sizeQR]{#2}
            \end{minipage}%
          };
      \fi
    \fi
  },
}

% ============================================================
%  BOX: MDFRAMED
% ============================================================
\usepackage{mdframed}

\newenvironment{cyanmargin}[1][]
  {\begin{mdframed}[topline=false,bottomline=false,rightline=false,
    linecolor=cyan,linewidth=.05cm,frametitlefont=\itshape]#1\ignorespaces}
  {\end{mdframed}}

\newenvironment{mdexercise}[1][]
  {\begin{mdframed}[topline=false,bottomline=false,rightline=false,
    linecolor=cyan,linewidth=.05cm,frametitlefont=\itshape]
   \textbf{Esercizio. }#1\ignorespaces}
  {\end{mdframed}}

\newenvironment{mdexam}[1][]
  {\begin{mdframed}[topline=false,bottomline=false,rightline=false,
    linecolor=magenta,linewidth=2,frametitlefont=\itshape]
   \textbf{Allenamento per la verifica. }#1\ignorespaces}
  {\end{mdframed}}

\newenvironment{mdproof}[1][]
  {\begin{mdframed}[topline=false,bottomline=false,rightline=false,
    linecolor=black,frametitle=Proof.,frametitlefont=\itshape]#1\ignorespaces}
  {\unskip\begin{flushright}\qedsymbol\end{flushright}\end{mdframed}}

\newenvironment{nestedproof}[1][\small]
  {\begin{mdframed}[topline=false,bottomline=false,rightline=false,
    linecolor=black,frametitle=Proof.,frametitlefont=\itshape]#1\ignorespaces}
  {\unskip\begin{flushright}\qedsymbol\end{flushright}\end{mdframed}}

% ============================================================
%  TEOREMI
% ============================================================
\newtheoremstyle{flemma}
  {3pt}{3pt}
  {\itshape\fontfamily{qpl}\selectfont}
  {1pt}{\bfseries}{ \space}{ }{}

\theoremstyle{definition}
\newtheorem*{definition}  {Definizione}
\newtheorem*{proposition} {Proposizione}
\newtheorem*{example}     {Esempio}
\newtheorem*{remark}      {Osservazione}
\newtheorem*{nota}        {Nota}
\newtheorem*{application} {Applicazione}

\theoremstyle{flemma}
\newtheorem*{property}  {Proprietà}
\newtheorem*{law}       {Legge}
\newtheorem*{lemma}     {Lemma}
\newtheorem*{theorem}   {Teorema}
\newtheorem*{corollary} {Corollario}
\newtheorem*{principle} {Principio}

\renewcommand\qedsymbol{\fontfamily{qpl}\selectfont\emph{Q.E.D.}}

% ============================================================
%  ESERCIZI E SOLUZIONI (ambienti classici)
% ============================================================
\usepackage{needspace}
\usepackage{xparse}
\usepackage{pifont}
\usepackage{environ}

\newcounter{exercise}
\renewcommand{\theexercise}{\arabic{exercise}}

% Stelline — nere, senza fbox
\newcommand{\onestar}  {\ding{72}\ding{73}\ding{73}\space}
\newcommand{\twostars} {\ding{72}\ding{72}\ding{73}\space}
\newcommand{\threestars}{\ding{72}\ding{72}\ding{72}\space}

\newcommand{\setexdifficulty}[1]{%
  \IfNoValueTF{#1}
    {\def\exdifficulty{\twostars}}
    {\def\exdifficulty{#1}}%
}

\NewDocumentEnvironment{exercise}{ o }{%
  \refstepcounter{exercise}\par\medskip
  \setexdifficulty{#1}\noindent
  \hypertarget{ex:\theexercise}{}\label{ex:\theexercise}%
  \textbf{Esercizio \theexercise\ }\exdifficulty\quad
  \hyperlink{sol:\theexercise}{\textit{(vai alla soluzione)}}\quad\ignorespaces
}{\par\medskip}

\NewDocumentEnvironment{question}{ o }{%
  \refstepcounter{exercise}\par\medskip
  \setexdifficulty{#1}\noindent
  \hypertarget{ex:\theexercise}{}\label{ex:\theexercise}%
  \textbf{Domanda di Teoria \theexercise\ }\exdifficulty\quad
  \hyperlink{sol:\theexercise}{\textit{(vai alla soluzione)}}\quad\ignorespaces
}{\par\medskip}

\NewDocumentEnvironment{exercisenosol}{ o }{%
  \refstepcounter{exercise}\par\medskip
  \setexdifficulty{#1}\noindent
  \hypertarget{ex:\theexercise}{}\label{ex:\theexercise}%
  \textbf{Esercizio \theexercise\ }\exdifficulty\quad\ignorespaces
}{\par\medskip}

\NewDocumentEnvironment{questionnosol}{ o }{%
  \refstepcounter{exercise}\par\medskip
  \setexdifficulty{#1}\noindent
  \hypertarget{ex:\theexercise}{}\label{ex:\theexercise}%
  \textbf{Domanda di teoria \theexercise\ }\exdifficulty\quad\ignorespaces
}{\par\medskip}

\NewDocumentEnvironment{expsolution}{ o }{%
  \Needspace{6\baselineskip}%
  \IfNoValueTF{#1}{\edef\resnum{\theexercise}}{\edef\resnum{#1}}%
  \par\bigskip
  \hypertarget{sol:\resnum}{}\label{sol:\resnum}%
  \noindent\textbf{Soluzione \hyperlink{ex:\resnum}{\resnum}}\quad
  \hyperlink{ex:\resnum}{\textit{(torna all'esercizio)}}\par\medskip
}{\par\medskip}

% solution: comportamento duale (inline grigio fuori da rigaesercizi,
%           accumulatore dentro rigaesercizi — vedi sotto)
\newenvironment{solution}
  {\par\noindent\hfill\small\color{gray}}
  {\par}

% ============================================================
%  MODALITÀ ESERCIZI (due colonne + testata ciano)
% ============================================================
\usepackage{multicol}
\usepackage{colortbl}

\newcommand{\exheadrule}[2]{%
  \setlength{\headheight}{20pt}%
  \fancyhead[RE]{%
    \colorbox{ciano}{%
      \parbox[c][14pt][c]{\textwidth}{%
        \color{white}\bfseries\sffamily\small #1\hfill}}}%
  \fancyhead[LO]{%
    \colorbox{ciano}{%
      \parbox[c][14pt][c]{\textwidth}{%
        \color{white}\bfseries\sffamily\small\hfill #2}}}%
  \fancyhead[RO,LE]{}%
  \renewcommand{\headrulewidth}{0pt}%
}

\newif\ifinsidemodalita
\insidemodalitafalse

\newenvironment{modalitaesercizi}{%
  \insidemodalitatrue
  \edef\exChapNum{\thechapter}%
  \edef\exChapName{\currentchaptername}%
  \renewcommand{\Needspace}[1]{}%
  \newgeometry{
    top=0cm, bottom=.8cm,
    inner=1.2cm, outer=1.2cm,
    includeheadfoot,
  }%
  \pagestyle{fancy}\fancyhf{}\fancyfoot[C]{\thepage}%
  \exheadrule%
    {CAPITOLO \exChapNum \enspace\textbar\enspace \exChapName}%
    {CAPITOLO \exChapNum \enspace\textbar\enspace ESERCIZI}%
  \setlength{\columnsep}{1cm}%
  \raggedcolumns
  \begin{multicols}{2}%
}{%
  \end{multicols}%
  \restoregeometry%
  \pagestyle{fancy}\fancyhf{}\fancyhead{}\fancyfoot[C]{\thepage}%
  \renewcommand{\headrulewidth}{0pt}%
  \insidemodalitafalse
}

\newcommand{\pausamulticol}{%
  \ifinsidemodalita\end{multicols}\par\vspace{2pt}\fi
}
\newcommand{\riprendimulticol}{%
  \ifinsidemodalita\par\vspace{2pt}\begin{multicols}{2}\fi
}

% ============================================================
%  AMBIENTE truthorfalse
%  Da aggiungere nel preambolo, dopo il caricamento di
%  tcolorbox (già presente) e enumitem (già presente).
%
%  Uso:
%    \begin{truthorfalse}
%      \tfitem{testo dell'affermazione}
%    \end{truthorfalse}
%
%  Le caselle V / F vengono inserite automaticamente
%  a destra di ogni voce.
% ============================================================

% Casellina grigia piccola (V oppure F).
% Si usa \newtcbox invece di \tcbox[...]{...} inline:
% \tcbox con opzioni esplicite dentro \newcommand causa
% conflitti nel parsing degli argomenti di tcolorbox.
% \newtcbox pre-registra le opzioni ed è privo di ambiguità.
\newtcbox{\tfvf}{%
  on line,
  arc=0.8pt, outer arc=0.8pt,
  boxrule=0.3pt, boxsep=0pt,
  left=2pt, right=2pt, top=0.8pt, bottom=0.8pt,
  colframe=gray!40, colback=gray!10,
  fontupper={\fontsize{6}{6}\selectfont\bfseries\color{gray!65}}%
}

% Singola voce: testo → caselle V F allineate a destra
\newcommand{\tfitem}[1]{%
  \item #1\nobreak\hfill\tfvf{V}\kern2pt\tfvf{F}%
}

% Lista vera/falsa
\newenvironment{truthorfalse}{%
  \begin{enumerate}[%
    label  = \textbf{\alph*.},%
    leftmargin = *,%
    itemsep = 3pt,%
    topsep  = 2pt,%
    parsep  = 0pt%
  ]%
}{%
  \end{enumerate}%
}
% ============================================================
%  Flag per raccolta soluzioni
% ============================================================
% ============================================================
%  *** SOSTITUIRE TUTTO da qui fino a fine \newenvironment{exrcell} ***
% ============================================================

% Flag per raccolta soluzioni
\newif\ifincollectmode  \incollectmodefalse
\newif\iffirstsolution  \firstsolutiontrue

% ============================================================
%  CORREZIONE BUG soluzioni duplicate con rigaesercizi{3}
%
%  Causa: il vecchio meccanismo con \newtoks memorizzava
%  \the\exrsoltemp come riferimento LAZY dentro \exrsoltoks.
%  Quando la terza soluzione sovrascriveva \exrsoltemp,
%  entrambi i riferimenti espandevano lo stesso valore:
%  risultato [sol1, sol3, sol3] invece di [sol1, sol2, sol3].
%
%  Fix: ogni soluzione viene copiata permanentemente in
%  \exrsol@1, \exrsol@2, \exrsol@3  via \let, che fotografa
%  \BODY al momento della chiamata.  Ridefinizioni successive
%  di \BODY non toccano le copie già salvate.
% ============================================================

\newcounter{exrsolcount}   % conta le soluzioni della riga corrente

\RenewEnviron{solution}{%
  \ifincollectmode
    \iffirstsolution
      \global\firstsolutionfalse
    \fi
    \stepcounter{exrsolcount}%
    % Copia permanente: \exrsol@N ottiene il significato corrente
    % di \BODY; eventuali ridefinizioni future di \BODY non la toccano
    \expandafter\global\expandafter\let
      \csname exrsol@\theexrsolcount\endcsname\BODY
  \else
    \par\noindent\hfill\small\color{gray}\BODY\par
  \fi
}

% ============================================================
%  Dimensioni e box ausiliari (invariati)
% ============================================================
\newsavebox{\exrfulltmpbox}
\newsavebox{\exrsolbox}
\newsavebox{\exrstarbox}
\newlength{\exrlabelwidth}
\newlength{\exrcellwidth}
\newlength{\exrsolcolwidth} \setlength{\exrsolcolwidth}{3.5cm}
\newlength{\exrcellsep}     \setlength{\exrcellsep}{6mm}
\newlength{\exrlabelgap}    \setlength{\exrlabelgap}{4mm}
\newlength{\exrtotalgap}
\newcounter{exrowcell}
\newcounter{exrowmaxcolsctr}

% ============================================================
%  Stampa le soluzioni accumulate come  [sol1; sol2; sol3]
%  Gestisce fino a 3 soluzioni per riga (massimo supportato).
% ============================================================
\newcommand{\exrprintsol}{%
  \savebox{\exrsolbox}{%
    \color{gray!70!black}%
    \ifnum\value{exrsolcount}>0\relax
      \csname exrsol@1\endcsname
      \ifnum\value{exrsolcount}>1\relax
        ;\space\csname exrsol@2\endcsname
        \ifnum\value{exrsolcount}>2\relax
          ;\space\csname exrsol@3\endcsname
        \fi
      \fi
    \fi
  }%
  \color{gray!70!black}%
  $\left[\rule[-\dp\exrsolbox]{0pt}%
    {\dimexpr\ht\exrsolbox+\dp\exrsolbox\relax}\right.$%
  \usebox{\exrsolbox}%
  $\left.\rule[-\dp\exrsolbox]{0pt}%
    {\dimexpr\ht\exrsolbox+\dp\exrsolbox\relax}\right]$%
}

\newenvironment{rigaesercizi}[2][\twostars]{%
  \setcounter{exrowcell}{0}%
  \setcounter{exrowmaxcolsctr}{#2}%
  \setcounter{exrsolcount}{0}%
  \global\firstsolutiontrue%
  \incollectmodetrue%
  \refstepcounter{exercise}%
  \setexdifficulty{#1}%
  \hypertarget{ex:\theexercise}{}\label{ex:\theexercise}%
  \sbox{\exrstarbox}{{\scriptsize\exdifficulty}}%
  \setlength{\exrlabelwidth}{\dimexpr\wd\exrstarbox + 6pt\relax}%
  \setlength{\exrtotalgap}{#2\exrcellsep}%
  \advance\exrtotalgap by -\exrcellsep%
  \ifnum#2=1\relax
    \setlength{\exrcellwidth}{\dimexpr
      \linewidth - \exrlabelwidth - \exrlabelgap\relax}%
  \else
    \setlength{\exrcellwidth}{\dimexpr
      ( \linewidth
        - \exrlabelwidth
        - \exrlabelgap
        - \exrsolcolwidth
        - \exrcellsep
        - \exrtotalgap
      ) / #2\relax}%
  \fi
  \par\nopagebreak\vspace{2pt}%
  % \vbox\bgroup\noindent\leavevmode
  \vbox\bgroup\noindent\leavevmode\parfillskip=0pt
  {\setlength{\fboxsep}{0pt}%
   \vtop{\offinterlineskip
     \hbox{\colorbox{intestRiassunto}{%
       \color{white}\bfseries\small
       \makebox[\exrlabelwidth][c]{\rule{0pt}{11pt}\theexercise}}}%
     \hbox{\makebox[\exrlabelwidth][c]{\scriptsize\exdifficulty}}%
   }}%
  \hspace{\exrlabelgap}\small\raggedright\parfillskip=0pt
}{%
  \iffirstsolution\else
    \unskip\nobreak\hfill\allowbreak\null\hfill\mbox{\exrprintsol}%
  \fi
  \par\egroup
  \vspace{4pt}%
  \incollectmodefalse%
}

\newenvironment{exrcell}{%
  \stepcounter{exrowcell}%
  \ifnum\value{exrowmaxcolsctr}=1\relax
    \ignorespaces
  \else
    \vtop\bgroup
      \hsize=\exrcellwidth
      \binoppenalty=10000
      \relpenalty=10000
      \small\raggedright\ignorespaces
  \fi
}{%
  \ifnum\value{exrowmaxcolsctr}=1\relax
    \unskip
  \else
    \egroup
    \ifnum\value{exrowcell}<\value{exrowmaxcolsctr}\relax
      \hspace{\exrcellsep}%
    \fi
  \fi
}

% ============================================================
%  *** FINE sezione da sostituire ***
% ============================================================
% ============================================================
%  SEZIONI INTERNE AL RIASSUNTO
% ============================================================
\newcommand{\rsection}[1]{%
  \par\vspace{0.55\baselineskip}%
  \begin{tcolorbox}[
    enhanced, breakable,
    colback=rsezColor, colframe=rsezColor,
    arc=2mm, boxrule=0pt,
    top=2mm, bottom=2mm, left=5mm, right=5mm,
  ]
    \color{white}\bfseries\sffamily\large #1%
  \end{tcolorbox}%
  \vspace{0.1\baselineskip}%
}

\newcommand{\rsubsection}[1]{%
  \par\vspace{0.35\baselineskip}%
  \noindent{\color{intestRiassunto}\bfseries\sffamily #1}%
}

\newcommand{\rkeyword}[1]{%
  \textcolor{intestRiassunto}{\bfseries #1}%
}

% ============================================================
%  AMBIENTE RIASSUNTO
% ============================================================
\newenvironment{riassunto}{%
  \clearpage
  \pagecolor{sfondoRiassunto}%
  \newgeometry{
    top    = 1.8cm,
    bottom = 1.8cm,
    inner  = 1.5cm,
    outer  = 1.5cm,
    includeheadfoot,
  }%
  \pagestyle{fancy}\fancyhf{}%
  \setlength{\headheight}{22pt}%
  \fancyhead[LE,RO]{%
    \colorbox{intestRiassunto}{%
      \parbox[c][16pt][c]{\textwidth}{%
        \color{white}\bfseries\sffamily\small
        \hspace{3mm}CAPITOLO~\thechapter
        \enspace\textbar\enspace
        \MakeUppercase{\currentchaptername}%
        \hfill RIASSUNTO\hspace{3mm}%
      }%
    }%
  }%
  \fancyhead[LO,RE]{}%
  \renewcommand{\headrulewidth}{0pt}%
  \fancyfoot[C]{\color{intestRiassunto}\thepage}%
  \par\vspace{0.6em}%
  \begin{center}
    {\color{intestRiassunto!50!black}\rule{0.18\textwidth}{1pt}}%
    \quad
    {\color{intestRiassunto}\fontsize{27}{31}\selectfont\bfseries\scshape Riassunto}%
    \quad
    {\color{intestRiassunto!50!black}\rule{0.18\textwidth}{1pt}}\\[0.35em]
    {\color{rsezColor!70!white}\fontsize{9}{11}\selectfont\sffamily
      Capitolo~\thechapter\ \textbullet\ \currentchaptername}%
  \end{center}%
  \vspace{0.7em}%
}{%
  \clearpage
  \pagecolor{white}%
  \restoregeometry%
  \pagestyle{fancy}%
  \fancyhf{}\fancyhead{}\fancyfoot[C]{\thepage}%
  \renewcommand{\headrulewidth}{0pt}%
}

% ============================================================
%  AMBIENTE VERIFICA SUBITO
% ============================================================
\newenvironment{verificasubito}[2][Verifica subito!]{%
  \def\vsncols{#2}%
  \renewcommand{\Needspace}[1]{}%
  \checkoddpage%
  \ifoddpage%
    \def\sforoSin{0cm}%
    \def\sforoDes{0cm}%
  \else%
    \def\sforoSin{0cm}%
    \def\sforoDes{0cm}%
  \fi%
  \begin{tcolorbox}[
    enhanced, breakable,
    colback=white, colframe=ciano,
    boxrule=1.5pt, arc=3mm,
    top=4mm, bottom=4mm, left=4mm, right=4mm,
    enlarge left by=\sforoSin,
    enlarge right by=\sforoDes,
    title={\large\bfseries\sffamily #1},
    coltitle=white,
    attach boxed title to top left={yshift=-2mm, xshift=4mm},
    boxed title style={
      empty,
      top=1.5mm, bottom=1.5mm, left=3mm, right=8mm,
      underlay={
        \fill[ciano, rounded corners=2mm]
          (frame.south west) -- (frame.north west) --
          ([xshift=-5mm]frame.north east) -- (frame.south east) -- cycle;
      }
    },
  ]%
  \setlength{\columnsep}{1cm}%
  \raggedcolumns%
  \ifnum\vsncols>1\relax
    \begin{multicols}{\vsncols}%
  \fi
}{%
  \ifnum\vsncols>1\relax
    \end{multicols}%
  \fi
  \end{tcolorbox}%
}

% ============================================================
%  QR NEL MARGINE
% ============================================================
\usepackage{qrcode}
\usepackage{marginnote}
\usepackage{marginfix}

\newcommand{\marginqr}[2][]{%
  \par\Needspace{6\baselineskip}\calcsizeQR
  \marginpar{\centering
    \begin{tcolorbox}[enhanced, boxrule=1.5pt, colframe=ciano, colback=white,
      arc=1mm, top=4pt, bottom=4pt, left=3pt, right=3pt,
      boxsep=0pt, width=\marginparwidth]
      \if\relax\detokenize{#1}\relax\else
        {\footnotesize\setlength{\parskip}{0pt}\setstretch{1.0}%
         \raggedright\sffamily #1\par\vspace{3pt}}
      \fi
      \centering\qrcode[height=\sizeQR]{#2}%
    \end{tcolorbox}}%
}


% ============================================================
%  AMBIENTE PROVATU — box nel margine, colori riassunto
% ============================================================
\newsavebox{\provatubox}
\newenvironment{provatu}[1][Prova tu!]{%
  \begin{lrbox}{\provatubox}%
    \begin{minipage}{\marginparwidth}%
      \begin{mdframed}[
        skipabove=0pt, skipbelow=0pt,
        linecolor=intestRiassunto,
        linewidth=1.2pt,
        backgroundcolor=sfondoRiassunto,
        frametitle={\color{white}\small\bfseries\sffamily\hspace{2pt}#1},
        roundcorner=4pt,
        frametitlebackgroundcolor=intestRiassunto,
        frametitlerule=false,
        frametitleaboveskip=2pt,
        frametitlebelowskip=2pt,
        innertopmargin=4pt,
        innerbottommargin=4pt,
        innerleftmargin=4pt,
        innerrightmargin=4pt,
      ]%
      \small\raggedright\ignorespaces
}{%
      \end{mdframed}%
    \end{minipage}%
  \end{lrbox}%
  \marginnote{\usebox{\provatubox}}%
  \ignorespaces
}{\ignorespacesafterend}

% ============================================================
%  TESTO E TIPOGRAFIA
% ============================================================
\usepackage{csquotes}
\usepackage[normalem]{ulem}
\usepackage{soul}
\usepackage{tipa}

\setulcolor{gray}
\renewcommand{\CancelColor}{\color{red}}

\newcommand{\schwa}{ə}
\newcommand{\Schwa}{Ə}
\newcommand{\lette}[1]{\enquote*{#1}}
\newcommand{\virgo}[1]{\enquote*{#1}}
\newcommand{\definizione}[1]{\textcolor{red}{\textbf{#1}}}
\newcommand{\citaz}[2]{\guillemotleft#1\guillemotright\begin{flushright}#2\end{flushright}}
\newcommand{\cit}[1]{\guillemotleft#1\guillemotright}
\newcommand{\boh}[1]{\textcolor{red}{#1}}
\newcommand{\bianco}[1]{\textcolor{white}{#1}}

% ============================================================
%  PACCHETTI GRAFICI E VARI
% ============================================================
\usepackage{graphicx}
\usepackage{array}
\usepackage{booktabs}
\usepackage{placeins}
\usepackage{titling}
\usepackage{float}
\usepackage{wrapfig}
\usepackage{sidecap}
\usepackage{subfig}
\usepackage{subcaption}
\usepackage{floatflt}
\usepackage{animate}
\usepackage{standalone}
\usepackage{listings}
\usepackage{footmisc}
\usepackage{xstring}
\usepackage{ifthen}
\usepackage[strict]{changepage}
\usepackage[margin=1cm]{caption}

% ============================================================
%  HYPERREF (sempre per ultimo)
% ============================================================
\usepackage[colorlinks]{hyperref}
\hypersetup{linkcolor=magenta}


\usepackage{lettrine}

\begin{document}
\chapter*{Funziona a meraviglia}
\begin{flushright}
    \textit{Wenke Husmann}\\
    {\small Internazionale, 22 gennaio 2026}
\end{flushright}

% Incolla il testo subito dopo \lettrine, a partire da "l cielo notturno..."
% (la "I" iniziale è il capolettera). Paragrafi separati da una riga vuota;
% i sottotitoli come \section*{Vantaggi evolutivi}, ecc.
\lettrine[lines=3]{I}{}l cielo notturno era attraversato da bagliori, come se degli alieni avessero acceso le luci nello spazio infinito. Era ottobre, e avevo svegliato mia figlia di buon’ora perché l’aurora boreale sarebbe stata visibile da Amburgo. Ci siamo vestite, siamo andate in un parco ed ecco il cielo illuminato da lunghe strisce infuocate. Rosso, viola, verde. I colori salivano e si trasformavano in giganteschi paesaggi di luce. Mia figlia guardava con gli occhi spalancati, mi ha preso la mano e non smetteva di esclamare “wow!”. Invece di stare al caldo nel letto come di solito a quell’ora, era su un prato gelato sotto l’enormità del cielo notturno, e ne era felicissima.

Mi è tornato in mente quando anche io, da bambina, ero piena di meraviglia: per la consistenza del fango tra le dita dei piedi, per le formiche che trasportavano pesi molto più grandi del loro, per la mia prima stella cadente. Davanti al cielo stellato sopra di me e agli occhi lucidi di mia figlia, mi sono chiesta dove sia finito il mio stupore. Dove sono le mie meraviglie?

Da quando, con l’illuminismo, gli scienziati hanno cominciato a indagare il mondo empiricamente, ogni sua meraviglia ha una spiegazione razionale. L’aurora boreale? Non è un messaggio mitologico, ma la collisione tra gli elettroni e gli atomi dell’atmosfera. I colori dei fiori? Non esprimono gli umori degli dèi, ma derivano dalla combinazione dei geni. I fuochi fatui nelle paludi? Non sono le anime dei defunti infelici, ma gas in autocombustione. All’inizio dell’era moderna il sociologo Max Weber scriveva che la meraviglia è un sentimento ingenuo, inferiore, provato solo “dalle anime semplici, dai matti e dai bambini”, o al massimo da “innamorati, poeti e devoti”. Ed è vero: nel mio mondo adulto, tutti i troll delle foreste si sono ridotti a semplici radici e il volto della Luna non è che la forma dei crateri sulla sua superficie. Appartengo a quella generazione di genitori secondo cui non è mai troppo presto per incoraggiare i figli a porsi domande e a sviluppare il pensiero critico.

È una cosa ragionevole. Però, mentre osservavo mia figlia sul prato gelato, provavo una certa nostalgia, o forse addirittura invidia. Non di saltellare in tondo come una bimba o di sentire in ogni conchiglia il suono di tutto l’universo. È che un mondo dove si può spiegare ogni magia è un mondo terribilmente triste. Posso imparare di nuovo la meraviglia? Esiste una versione adulta dell’incanto?

Prima di tutto è necessario chiarire i concetti. La meraviglia è quella sensazione di rapimento che ci coglie di fronte a qualcosa di nuovo, inaspettato, magnifico. Diversamente dallo spavento, non suscita la paura, semmai un timore reverenziale. Alla meraviglia si associano entusiasmo e ammirazione, che vanno oltre il semplice stupore. È un’emozione complessa e universale, a cui reagiamo ovunque nello stesso modo: piegando la testa all’indietro, spalancando occhi e braccia ed esclamando qualcosa tipo “wow!”.

Ricordo ancora quella sensazione. Quando ero bambina, era intensa. Come la paura, la felicità o la rabbia. E le emozioni intense hanno sempre avuto una funzione importante nell’evoluzione umana. La paura, per esempio, ci proteggeva dai predatori; la gioia ci spinge a cercare un partner e a riprodurci; la rabbia ci aiuta a superare gli ostacoli.

\subsection*{Vantaggi evolutivi}

Quindi anche la meraviglia deve avere una funzione. Ma quale? Per molto tempo gli scienziati si sono occupati più che altro delle emozioni che causano problemi: la rabbia, la tristezza, l’ansia, perché comprenderle ci aiuta a gestirle meglio. Solo in seguito hanno cominciato a indagare anche i sentimenti positivi come la sorpresa, la gioia, l’amore. E da circa vent’anni anche la meraviglia.

All’università della California a Berkeley, gli scienziati hanno studiato in che modo le varie emozioni positive influiscono su alcuni indici di infiammazione nel sangue. Hanno scoperto che ad abbassare certi valori è soprattutto la meraviglia. Altri studi hanno prodotto risultati simili per la pressione sanguigna e i livelli di cortisolo: tra le emozioni piacevoli è sempre la meraviglia a distinguersi per i suoi effetti positivi. Gli scienziati suppongono che potrebbe contribuire indirettamente anche ad alleviare il dolore, a migliorare la qualità del sonno, a evitare ansia e depressione. Quelle ricerche hanno posto le basi per una vera e propria scienza. Eftychia Stamkou è una giovane psicologa dell’università di Amsterdam che da tre anni studia la meraviglia nei bambini. In uno dei suoi esperimenti ha mostrato a tre gruppi di bambini tra gli otto e i tredici anni dei video che provocavano meraviglia, allegria o nessuna delle due emozioni. I risultati sono stati sorprendenti.


Stamkou ha misurato indirettamente l’attività del nervo vago, che collega il cervello al cuore, ai polmoni, all’intestino e ad altri organi. Come un istruttore di yoga interno, fa sì che dopo la fatica o lo stress torniamo a riposarci: il battito rallenta, il corpo si rilassa. Ci sentiamo bene. È la stessa reazione rilevata tra i bambini che hanno provato meraviglia nell’esperimento: erano più rilassati rispetto ai bambini degli altri due gruppi.

Stamkou ha registrato anche un altro effetto. Come premio, al termine dell’esperimento tutti i partecipanti potevano scegliere il loro dolce preferito. In seguito gli si chiedeva se preferivano tenerlo oppure donarlo a dei bambini rifugiati. Nel gruppo in cui non erano state stimolate particolari emozioni la maggior parte dei bambini ha tenuto il dolce per sé. Lo stesso è successo con i bambini che avevano provato allegria. Tra quelli che avevano provato meraviglia, invece, circa due terzi hanno donato il dolcetto, quasi il doppio rispetto agli altri due gruppi. “La meraviglia non solo li fa sentire bene”, spiega Stamkou, “ma li rende più generosi e può contribuire, attraverso piccoli gesti di gentilezza, a rendere il mondo ogni giorno migliore”. Com’è possibile?

A Berkeley, Dacher Keltner studia in che modo le emozioni influenzano la nostra comprensione del mondo e il nostro comportamento. È interessato soprattutto ai sentimenti positivi. Ha creato un centro interdisciplinare di ricerca e formazione, in cui benessere e felicità sono studiati dal punto di vista scientifico. Anche lui è rimasto colpito dalla meraviglia, a cui ha dedicato un intero libro (Wow, Il saggiatore 2023). “La meraviglia offre agli esseri umani vantaggi evolutivi decisivi”, spiega.


\subsection*{Piccoli e connessi}

Per scoprire quali, Keltner ha chiesto ad alcuni studenti di guardare per due minuti alcuni alberi di eucalipto alti sessanta metri, mentre un gruppo di controllo doveva guardare la facciata di un austero edificio universitario, della stessa altezza. Quando ai partecipanti è stato chiesto quanto volessero ricevere come rimborso spese, gli studenti che avevano ammirato gli eucalipti hanno chiesto una cifra minore rispetto a quelli del gruppo di controllo. E quando uno dei ricercatori ha fatto cadere a terra intenzionalmente libri e penne, gli studenti che avevano provato meraviglia hanno mostrato più solerzia nell’aiutarlo a raccoglierli rispetto agli altri. Gli adulti dell’esperimento di Keltner hanno reagito come i bambini dell’esperimento di Stamkou: anche loro si sono mostrati più generosi e premurosi. Un vantaggio evolutivo per gli umani, la cui sopravvivenza è legata strettamente alla vita in comunità.

Keltner definisce la meraviglia così: “È una sensazione che proviamo in presenza di qualcosa di grande, che va oltre la nostra comprensione del mondo”. Non si riferisce tanto alle dimensioni fisiche, quanto piuttosto a qualcosa di grandioso. Si può provare meraviglia anche di fronte alla perfezione di un fiocco di neve. La cosa importante è ciò che ne deriva: la sensazione di piccolezza del proprio io. L’essere umano diventa più umile.


Una collega di Keltner, la psicologa sociale Yang Bai, ha offerto una prova convincente al riguardo. Su una piattaforma che offre una vista spettacolare sul Grand canyon, ha chiesto a dei turisti di disegnare uno schizzo di se stessi. Come metro di paragone ha proposto lo stesso compito ai turisti che passeggiavano lungo il vivace lungomare di San Francisco. Risultato: negli schizzi realizzati sul Grand canyon le figure avevano dimensioni di un terzo più piccole rispetto a quelle disegnate a San Francisco. Dalle interviste alle persone che avevano provato meraviglia sul Grand canyon è risultato, inoltre, che si percepivano non tanto come individui ma come parte di un gruppo. Non è un caso che d’istinto, davanti all’aurora boreale, mia figlia e io ci siamo prese la mano.

La differenza fondamentale rispetto alle altre emozioni positive è proprio questa: chi la prova entra in connessione con il mondo. Mentre l’orgoglio o la soddisfazione, per quanto assolutamente piacevoli, sono rivolte all’interno, al proprio sé, davanti a qualcosa di grandioso ci si sente proiettati verso l’esterno, verso la comunità. In questo senso la meraviglia somiglia di più al divertimento e alla gioia. “Anche questi sentimenti ci mettono in relazione con gli altri”, spiega Keltner, “e sono importanti per le amicizie: si ride, si scherza e insieme ci si tira su. Ma è con la meraviglia che entra davvero in ballo la collettività”. La meraviglia ci rende più cooperativi e meno narcisisti. Come se per un momento l’ego fosse ridimensionato.

Per questo Keltner paragona l’effetto della meraviglia a quello delle droghe psichedeliche, anch’esse capaci di annullare i confini dell’io. Quando si è in un trip ci si sente connessi con tutto e con tutti. Il cervello attiva processi analoghi. Certo, nel caso della meraviglia in misura minore, per meno tempo e senza effetti collaterali, precisa Keltner. “Ma sembra che i racconti dei viaggiatori rapiti dal Grand canyon e quelli di chi ha avuto un’esperienza psichedelica siano molto simili”. La voce dell’ego si fa più sommessa, i suoni del mondo più forti e colorati, ci si sente in connessione con una sorta di potere superiore. Ci si dimentica un po’ di sé e ci si apre di più agli altri, alla grandezza del tutto.


Le sostanze allucinogene e la meraviglia riducono l’attività delle aree cerebrali che servono a mettere in guardia dai pericoli, l’amigdala in particolare. Per questo i medici curano alcuni pazienti affetti da disturbo post-traumatico da stress con piccole dosi di queste sostanze. “Chi prova meraviglia si sente meno oppresso dall’ansia, il mondo gli sembra più bello”, dice Keltner.

È esattamente ciò di cui ho bisogno quando la mattina mi sveglio e leggo le notizie dal mondo. Chiedo a Keltner se la capacità di meravigliarsi aiuta a gestire meglio situazioni complesse o stressanti. Sì, risponde lui. Perché il cambiamento di percezione ci porta a essere più curiosi verso le novità. Ed è proprio questo il secondo vantaggio evolutivo: chi prova meraviglia vuole scoprire di più. E impara più facilmente.

\subsection*{Lasciarsi trasportare}

Ormai mi è chiaro che entusiasmarsi per qualcosa di straordinario è come un superpotere. Rende più sani, più socievoli e potenzialmente più intelligenti. Ma non sono ancora sicura di poter riacquistare questo potere. Ci si può esercitare? E se sì, come? In fondo non sono più una bambina che scopre ogni giorno qualcosa di grandioso solo perché per lei tutto è nuovo. Ma non voglio nemmeno aspettare il prossimo fenomeno celeste.

“In linea di principio, gli adulti si meravigliano per le stesse cose dei bambini”, sostiene Stamkou. Cose che troviamo ovunque, ma soprattutto nella natura: le ali di farfalla, i cristalli di ghiaccio, una cascata, il cielo stellato. Il segreto è guardare con attenzione e lasciarsi trasportare. Durante una passeggiata o semplicemente mentre si va a lavorare, fermarsi a guardare un fiore delicato, uno scorcio, la chioma di un albero, e per un attimo lasciarsi sopraffare: che meraviglia!

Ma forse io sono troppo pragmatica. Non mi ci vedo molto ad ammirare gli alberi. Mi ritrovo di più nella passione personale di Stamkou, che nella sua ricerca si è concentrata sugli effetti dell’arte e della finzione.

Keltner paragona l’effetto della meraviglia alle droghe psichedeliche

“Nella musica, nei film, nella letteratura o nelle arti visive, si tratta sempre di mettere in discussione ciò che è familiare, cercare l’ignoto e lasciarsi stupire”, dice. L’arte è la porta magica per accedere a un altro modo di vedere. È questo che ci affascina e ci diverte. Uno spettacolo a teatro, l’entusiasmo per un concerto o le lacrime che verso guardando un film sono tutte esperienze che allenano alla meraviglia.

E poiché la meraviglia risveglia effettivamente la curiosità e la sete di sapere, fin dall’antichità è al centro delle teorie sulla conoscenza. La profonda meraviglia, spiegava Platone, è l’approccio fondamentale del filosofo che mira a cercare la verità nascosta dietro i fenomeni quotidiani. Nel Teeteto fa dire a Socrate che la filosofia comincia con la meraviglia. A Socrate è attribuita anche la frase “so di non sapere”. E chi prova meraviglia sa perlomeno di non sapere tutto. Un buon inizio per imparare cose nuove.

Di questo impulso all’apprendimento si servono anche oratori e poeti, che nei loro discorsi e nei loro testi giocano con gli effetti di straniamento e sorpresa. È lo stesso motivo per cui sono state costruite le piramidi e le cattedrali.

C’è di più: attraverso l’osservazione contemplativa e la percezione consapevole prepariamo la mente ad acquisire conoscenze razionali. Esiste una quantità di storie su illuminazioni raggiunte durante un’osservazione stupita: Isaac Newton sotto il melo, Cartesio che guarda un arcobaleno. E nella sua Autobiografia scientifica il fisico tedesco Max Planck scrisse: “La meraviglia per i miracoli della natura è all’origine di ogni pensiero scientifico”.

\subsection*{Dietro l’infinito}

Gli esseri umani provano meraviglia quando incontrano qualcosa che non rientra nella loro comprensione del mondo e costringe il cervello a sciogliere la contraddizione e a cambiare prospettiva. Gli psicologi parlano di ristrutturazione cognitiva e adattamento. Mentre elaborano le nuove informazioni, i neuroni trasformano la struttura del cervello.

Questo è forse il terzo grande vantaggio: la meraviglia aiuta ad affrontare la vita in maniera più giocosa. Gli scienziati parlano spesso della meraviglia ai confini della capacità di conoscenza umana. Per quanto la mente umana possa crescere, non potrà mai comprendere tutto. Il fisico danese Niels Bohr disse: “Chi non è rimasto scioccato dalla meccanica quantistica non l’ha capita”. Perché con la sua teoria della materia Bohr ha introdotto nella fisica qualcosa di sorprendente: il caso. Che, diversamente da Albert Einstein, Bohr vedeva non come un segno dell’incompletezza di questa teoria, ma come inevitabile e fondamentale. Questo ha scatenato discussioni accese sulla natura della realtà: il mondo è reale indipendentemente dal fatto che noi lo osserviamo? Esistono dei limiti a ciò che possiamo scoprire?

“Cosa c’è dietro quelle luci?”, mi ha chiesto mia figlia guardando l’aurora boreale. “L’infinito”, ho risposto. “E dopo?”, ha chiesto ancora. Abbiamo parlato un po’ di cosa potesse esserci: del nulla che circonda tutto. Ma non eravamo soddisfatte delle nostre ipotesi e abbiamo ripreso a guardare in su, in silenzio. Alla fine, mia figlia ha detto: “Dietro l’infinito c’è quello che non possiamo immaginare”.

Magari la magia della meraviglia sta proprio in questo: ai confini della nostra comprensione, dove comincia l’inconcepibile, cerchiamo uno spiraglio da cui scorgere l’essenza della realtà.

\hrulefill

Puoi vedere l'articolo per completo scannerizzando il seguente codice QR:

\begin{center}
    \qrcode{https://www.internazionale.it/magazine/wenke-husmann/2026/01/22/funziona-a-meraviglia}
\end{center}


\section*{La tragedia del Vajont}
Puoi vedere lo spettacolo teatrale di Marco Paolini per intero qui:

\begin{center}
    \qrcode{https://www.youtube.com/watch?v=q7CarpRflzs}
\end{center}

\end{document}
